\chapter*{Tóm tắt}

Các hệ thống cảnh báo khí gas truyền thống hoạt động theo nguyên tắc \textit{ngưỡng tĩnh} (static threshold): chỉ phát cảnh báo khi nồng độ khí đã vượt mức nguy hiểm – đúng lúc mà người dùng cũng có thể nhận ra bằng khứu giác. Phương pháp này về bản chất là phản ứng thụ động, để lại ít thời gian can thiệp và dẫn đến tỷ lệ cảnh báo sai cao.

Khóa luận này trình bày việc thiết kế và triển khai một \textbf{nền tảng IoT 4 lớp} cho phát hiện và ứng phó rò rỉ khí gas, trong đó đóng góp khoa học chính là thay thế cơ chế cảnh báo thụ động bằng hệ thống trí tuệ nhân tạo hai tầng:

\begin{enumerate}
  \item \textbf{Mô hình dự báo LSTM} học từ dữ liệu vật lý tổng hợp của bộ mô phỏng và tính toán xác suất $P(\text{nồng độ khí} > 1000\,\text{ppm trong 5 phút tới})$ theo thời gian thực, cảnh báo \textit{trước khi} khí đạt mức nguy hiểm.

  \item \textbf{Bộ điều khiển học tăng cường PPO} (Proximal Policy Optimization) được huấn luyện trong môi trường Gymnasium tích hợp vật lý của bộ mô phỏng, mỗi giây lựa chọn một trong bốn hành động: \texttt{NO\_OP / ALERT\_USER / FAN\_ON / CLOSE\_VALVE}. Chính sách học được ưu tiên ngăn chặn rò rỉ ở mức chi phí thấp nhất, đồng thời phạt cả trường hợp bỏ lỡ lẫn cảnh báo sai.
\end{enumerate}

Hệ thống được triển khai theo kiến trúc 4 lớp: (1) \textbf{Lớp thiết bị} – ESP32 hoặc bộ mô phỏng phần mềm phát dữ liệu qua MQTT; (2) \textbf{Lớp truyền thông} – Mosquitto broker và dịch vụ cầu nối MQTT-Kafka; (3) \textbf{Lớp xử lý} – Apache Spark Structured Streaming chạy song song mô hình LSTM và bộ điều khiển RL, ghi kết quả vào InfluxDB và PostgreSQL; (4) \textbf{Lớp ứng dụng} – API Node.js/TypeScript, dashboard web thời gian thực, Grafana và chatbot Telegram cho phép nhận cảnh báo và điều khiển từ xa.

Để đánh giá khoa học, một kịch bản benchmark xác định (cùng seed ngẫu nhiên) được áp dụng để so sánh ba bộ điều khiển \texttt{THRESHOLD}, \texttt{FORECASTER} và \texttt{RL} trên các chỉ số: thời gian cảnh báo trước (\textit{lead time}), tỷ lệ bỏ sót, số cảnh báo sai mỗi giờ và nồng độ đỉnh trung bình trong pha rò rỉ. Kết quả cho thấy bộ điều khiển RL vượt trội so với ngưỡng tĩnh về thời gian cảnh báo sớm và khả năng giảm thiểu nồng độ đỉnh, trong khi bộ dự báo LSTM cải thiện đáng kể tỷ lệ bỏ sót so với phương pháp ngưỡng đơn thuần.

Toàn bộ hệ thống được đóng gói bằng Docker Compose, đảm bảo khả năng tái tạo và triển khai dễ dàng.

\bigskip
\noindent\textbf{Từ khóa:} Phát hiện rò rỉ khí gas, IoT, LSTM, Học tăng cường, PPO, Kafka, Apache Spark, Telegram, Docker.
